%==============================================================================
%  03-definitions.tex — Definitions & Conceptual Foundations
%==============================================================================
\strategyPart{03 \textbar\ Definitions}{Conceptual Foundations}

Clear definitions prevent the strategy from collapsing into technology
shopping. This section defines the AI systems organizations deploy and, more
importantly, the four organizational states those systems produce.

\section{Definitions of AI systems}

\begin{itemize}
  \item \textbf{Artificial intelligence} — systems performing activities
        associated with human intelligence: prediction, classification,
        language understanding, planning, generation and decision support
        [OECD AI Principles; EU AI Act; Stanford AI Index].
  \item \textbf{Generative AI} — systems that generate text, images, code,
        audio, video and structured data from learned patterns
        [Stanford AI Index; NIST Generative AI Profile].
  \item \textbf{AI model} — generates predictions, classifications,
        recommendations or content (LLMs, predictive, recommendation,
        classification and vision models).
  \item \textbf{AI assistant} — supports a person but normally does not act
        independently (writing, research, product-management and BI
        assistants).
  \item \textbf{AI copilot} — works inside a business process and suggests
        actions while the employee remains responsible (sales, coding,
        financial-analysis and service copilots).
  \item \textbf{AI agent} — receives an objective, determines the steps,
        uses tools and may execute actions with limited human intervention.
  \item \textbf{AI workflow} — a fixed, engineer-controlled sequence of AI
        steps; use it when the process is structured and rules are stable
        [Anthropic].
  \item \textbf{Multi-agent system} — several specialized agents cooperate
        (planning, retrieval, analysis, validation, reporting, governance).
  \item \textbf{AI-orchestrated organization} — agents coordinate
        information, tasks, decisions and workflows across teams and systems.
\end{itemize}

\section{The agent vocabulary at a glance}

\begin{center}
  \fitwidth{\begin{tabular}{>{\raggedright\arraybackslash}p{2.2cm}>{\raggedright\arraybackslash}p{3.3cm}>{\raggedright\arraybackslash}p{1.9cm}>{\raggedright\arraybackslash}p{2.4cm}>{\raggedright\arraybackslash}p{2.5cm}}
    \rowcolor{inNavy}
    \hcell{System} & \hcell{Main role} & \hcell{Autonomy} & \hcell{Human role} & \hcell{Main risk} \\
    \textbf{Assistant} & Provides information & Low & User & Incorrect info \\
    \textbf{Copilot} & Suggests in a process & Low--medium & Reviewer & Overreliance \\
    \textbf{Workflow} & Executes fixed sequence & Fixed & Designer & Brittleness \\
    \textbf{Agent} & Completes bounded tasks & Medium--high & Supervisor & Unauthorized action \\
    \textbf{Multi-agent} & Coordinates specialists & Variable & Governor & Cascading failure \\
  \end{tabular}}
\end{center}

The central design question is not ``can an agent replace a person'' but
\textbf{which parts of a job AI can perform safely, which parts it should
assist with, and which responsibilities must remain under human judgment and
accountability} [Anthropic; NIST AI RMF; Microsoft agent framework].
Anthropic's agent research recommends \textbf{selecting the simplest
architecture that reliably solves the business problem} — using predictable
workflows when the process is structured, and agents only when the task
requires dynamic, flexible interaction.

\section{Adoption, enablement, transformation and AI-native}

\begin{center}
  \renewcommand{\arraystretch}{1.5}
  \rowcolors{2}{inSnow}{white}
  \fitwidth{\begin{tabular}{>{\raggedright\arraybackslash}p{3.4cm}>{\raggedright\arraybackslash}p{4.6cm}>{\raggedright\arraybackslash}p{4.6cm}}
    \rowcolor{inNavy}
    \hcell{Concept} & \hcell{Description} & \hcell{Organizational signal} \\
    \textbf{AI adoption} & Employees or teams use individual AI tools. & Usage statistics; no process change. \\
    \textbf{AI enablement} & AI is integrated into selected business processes. & Process redesign in a few functions. \\
    \textbf{AI transformation} & Work, structures, capabilities and governance are redesigned around AI. & Enterprise-wide redesign, new roles, P\&L impact. \\
    \textbf{AI-native} & AI is the foundation of products, operations and decision-making from day one. & Business model designed around AI. \\
  \end{tabular}}
\end{center}

The distinction matters because the statistics are asymmetric: 88\% of
organizations have \textit{adopted} AI, but only 7\% have \textit{scaled} it
[McKinsey State of AI 2025]. Most organizations that call themselves ``AI
companies'' are, by these definitions, still in adoption or early
enablement. The remainder of this document is a route-map between those two
states.
